\chapter{Continuity and concentration in cycle families}\label{sec:sweepout-concentration}
The admissibility condition in Definition \ref{def:no-concentration} is uniform in the parameter and the center of the ball. We prove sufficient conditions and construct a flat-continuous family that fails it. Except for the final counterexample about dimension zero, $M^{n+1}$ is smooth, closed and connected, $n\geq1$, and cycles are integral mod-two boundaries as in Definition \ref{def:minmax-cycles}. The primary source is \cite{LMN}, Section 2.6 and Lemma 2.7, pp.~941--942.

\begin{lemma}[Compact families of nonatomic measures have uniform small-ball control]\label{lem:compact-nonatomic-measures}
Let $Y$ be a compact metric space and let $\mathcal K$ be a weakly compact family of finite nonnegative Radon measures on $Y$. Suppose every member of $\mathcal K$ has no atoms. Then
\[
 \lim_{r\downarrow0}\sup_{\mu\in\mathcal K,\,y\in Y}\mu(B_r(y))=0.
\]
The total masses are uniformly bounded.
\end{lemma}
\begin{proof}
Evaluation against the constant function one is continuous, so weak compactness bounds the total masses. On a compact metric space the weak topology on bounded sets of nonnegative measures is metrizable, by testing a countable dense subset of $C(Y)$. If the assertion failed, there would be $\varepsilon>0$, $r_j\downarrow0$, $y_j\in Y$, and $\mu_j\in\mathcal K$ with $\mu_j(B_{r_j}(y_j))\geq\varepsilon$. Pass to a subsequence with $y_j\to y$ and $\mu_j\rightharpoonup\mu\in\mathcal K$. For every $\rho>0$, the small balls eventually lie in $\overline B_\rho(y)$. The closed-set inequality for weak convergence gives
\[
 \varepsilon\leq\limsup_j\mu_j(\overline B_\rho(y))\leq\mu(\overline B_\rho(y)).
\]
Let $\rho\downarrow0$. Continuity from above for the finite measure gives $\mu(\{y\})\geq\varepsilon$, a contradiction. The closed-set inequality used here is the Portmanteau theorem, an external elementary weak-convergence fact. It is essential that the limiting measure belongs to the nonatomic family.
\end{proof}

\begin{proposition}[Compact F-continuous cycle families are admissible analytically]\label{prop:f-continuity-no-concentration}
Let $X$ be a compact metric parameter space and let $\Phi:X\to\mathcal Z_n^0(M;\mathbb Z_2)$ be $\mathbf F$-continuous, with $n\geq1$. Then $\Phi$ has bounded maximal mass and no concentration of mass. No cohomological detection is implied by this assertion. The topology is that of Definition \ref{def:minmax-cycles}, from \cite{MNIndex}, Section 2.
\end{proposition}
\begin{proof}
The associated varifold map $x\mapsto|\Phi(x)|$ is weakly continuous by the definition of $\mathbf F$. Testing a continuous function on $M$ as a function independent of the tangent plane gives weak continuity of the weight map $x\mapsto\|\Phi(x)\|$. Its image is a compact family of finite measures. Each weight is that of a rectifiable $n$-cycle; since $n\geq1$, it gives zero mass to every singleton. Lemma \ref{lem:compact-nonatomic-measures} gives both conclusions. In particular the varifold term controls a feature that the flat distance alone need not control.
\end{proof}

\begin{proposition}[Mass-continuous families have no concentration]\label{prop:mass-continuity-no-concentration}
Let $X$ be compact and let $\Phi:X\to\mathcal Z_n^0(M;\mathbb Z_2)$ be continuous in the mass norm, with $n\geq1$. Then its maximal mass is bounded and it has no concentration of mass. This is the closed-manifold, positive-dimensional case of \cite{LMN}, Lemma 2.7.
\end{proposition}
\begin{proof}
The image is compact in the mass metric. Given $\varepsilon>0$, choose finitely many cycles $T_1,\ldots,T_N$ so every $\Phi(x)$ is within mass distance $\varepsilon/2$ of one of them. Each fixed weight is finite and nonatomic, so Lemma \ref{lem:compact-nonatomic-measures}, applied to this finite family, gives $r>0$ such that $\|T_i\|(B_r(y))<\varepsilon/2$ for all $i,y$. The local mass triangle inequality gives
\[
 \|\Phi(x)\|(B_r(y))
 \leq\|T_i\|(B_r(y))+\mathbf M(\Phi(x)-T_i)<\varepsilon.
\]
Thus the small-ball estimate is uniform. The global triangle inequality with the same finite covering bounds all masses. Also $\mathcal F(T-S)\leq\mathbf M(T-S)$, so the family is flat-continuous. This proof does not identify every flat-continuous family with a mass-continuous one; the approximation theorems of \cite{LMN}, Section 2.9, are separate inputs.
\end{proof}

\begin{example}[Many small spheres concentrate while converging flatly to zero]\label{ex:flat-cycle-concentration}
On the flat unit torus $M=\mathbb R^{n+1}/\mathbb Z^{n+1}$, $n\geq1$, there are smooth mod-two boundary cycles $T_j$ and constants $c,C>0$ such that
\[
 \mathcal F(T_j)\longrightarrow0,\qquad c\leq\mathbf M(T_j)\leq C,
 \qquad \operatorname{spt}T_j\subset B_{1/(5j)}(0).
\]
Every individual weight is nonatomic, but the sequence has no uniform small-ball control.
\end{example}
\begin{proof}
The ball of radius $1/5$ is Euclidean. Put $r_j=1/(10j)$, $N_j=j^n$, and
\[
 \rho_{j,k}=r_j\left(1+\frac{k}{2N_j}\right),\qquad
 T_j=\sum_{k=1}^{N_j}\partial[B_{\rho_{j,k}}(0)]\pmod2.
\]
The spheres are distinct and disjoint, so their masses add. If $\sigma_n=\mathcal H^n(S^n)$, then
\[
 \sigma_n10^{-n}\leq\mathbf M(T_j)
 \leq\sigma_n(3/2)^n10^{-n}.
\]
The top-dimensional mod-two chain $U_j=\sum_k[B_{\rho_{j,k}}(0)]$ has boundary $T_j$. Its multiplicity is zero or one and its support lies in $B_{3/(20j)}(0)$, so
$\mathcal F(T_j)\leq\mathbf M(U_j)\leq\omega_{n+1}(3/(20j))^{n+1}\to0$.
The supports also lie in the larger displayed ball. At radii $1/(5j)$ all of the positive lower mass is present, proving the failure of a uniform small-ball estimate. Each slice is a finite union of smooth positive-dimensional spheres and therefore has no atoms.
\end{proof}

\begin{example}[A flat-continuous interval family with concentration]\label{ex:flat-continuous-concentrating-family}
There is a flat-continuous map $\Phi:[0,1]\to\mathcal Z_n^0(M;\mathbb Z_2)$ on the flat torus of Example \ref{ex:flat-cycle-concentration}, with bounded maximal mass, that fails Definition \ref{def:no-concentration}. Thus even a finite one-dimensional parameter complex does not remove this issue.
\end{example}
\begin{proof}
For $I_j=[2^{-(j+1)},2^{-j}]$, choose a continuous triangular function $a_j:I_j\to[0,1]$ that is zero at both endpoints and one at the midpoint. On $I_j$ let $\Phi(t)$ be the sum of the concentric spheres from Example \ref{ex:flat-cycle-concentration}, with every radius multiplied by $a_j(t)$. Interpret radius zero as the zero cycle. Set $\Phi=0$ on $[1/2,1]$ and at zero.
For fixed $j$, two nearby parameters are flat-close: their difference is filled by the sum, modulo two, of the annular regions between corresponding spheres, whose total volume tends to zero. At the endpoints the spheres collapse flatly to zero. Uniformly for $t\in I_j$, a mod-two filling is supported in $B_{3/(20j)}(0)$ and has mass at most its volume. This proves continuity at the accumulating endpoint $t=0$ as well. Scaling multiplies each sphere's area by $a_j(t)^n$, so all masses remain bounded. At the midpoints the family equals $T_j$, hence fails uniform no concentration by the preceding example. It is not asserted to be a sweepout: indeed the interval has no nonzero first cohomology class.
\end{proof}

\begin{example}[Why positive cycle dimension is required]\label{ex:zero-cycle-concentration}
Mass continuity alone does not imply no concentration for zero-dimensional cycles. A constant family whose value is the mod-two boundary of a nontrivial short arc in a circle is a counterexample to such a zero-dimensional extension of Proposition \ref{prop:mass-continuity-no-concentration}.
\end{example}
\begin{proof}
The boundary consists of its two distinct endpoints, each with unit mass. The constant family is continuous in every cycle topology, but every ball centered at one endpoint has mass at least one. Thus the small-ball supremum never tends to zero. The nonatomic step in the positive-dimensional proof is exactly what fails. All volume-spectrum statements in this project retain $n\geq1$.
\end{proof}
